\documentclass[12pt]{article}
\usepackage[utf8]{inputenc}
\usepackage[T1]{fontenc}
\usepackage{hyperref}
\title{Bukpe Language Resources for K-12 Curriculum Development}
\author{Compiled by Claude Sonnet 5 (claude-sonnet-5) for Andr\'as Kornai}
\date{}
\begin{document}
\maketitle
\section{Introduction}
Bukpe (also called Ukpe-Bayobiri, ISO 639-3 ukp) is a language of Cross River State, Nigeria. It is reported in Obudu and Ikom Local Government Areas and belongs to the Bendi branch of Benue--Congo. Sources give about 21,500 speakers (2013) and, for the people group, 28,000 to 29,000. Bayobiri is listed as a dialect.
\section{Primary Educational Resources}
None were found in search results. Joshua Project lists the language as not known to be taught in schools.
\section{Linguistic and Reference Materials}
\begin{itemize}
\item Wikipedia, Ukpe-Bayobiri language: \url{https://en.wikipedia.org/wiki/Ukpe-Bayobiri_language}
\end{itemize}
No dictionary, grammar or orthography was found.
\section{Religious and Community Resources}
\begin{itemize}
\item Joshua Project, Ukpe-Bayobiri in Nigeria; the scripture status is listed as translation started: \url{https://joshuaproject.net/people_groups/15695}
\item Joshua Project language page for ukp: \url{https://www.joshuaproject.net/index.php/languages/ukp}
\end{itemize}
\section{Digital and Audio Resources}
\begin{itemize}
\item Global Recordings Network, Bukpe language page: \url{https://www.globalrecordings.net/en/language/ukp}
\item Global Recordings Network, Ukpe-Bayogin: Ukpe language page: \url{https://www.globalrecordings.net/en/language/2285}
\end{itemize}
\section{Curriculum Development Context}
Little exists. The search found only reference and mission-oriented pages, with audio recordings the likeliest starting point for native-language content. A curriculum would need community orthography and materials development from the ground up.
\end{document}